\documentclass[a4,11pt]{aleph-notas}
% Se puede ver la documentación aquí:
% https://github.com/alephsub0/LaTeX_aleph-notas

% -- Paquetes adicionales
\usepackage{enumitem}
\usepackage{url}
\usepackage{array}
\usepackage{booktabs}
\usepackage{longtable}
\hypersetup{
    urlcolor=blue,
    linkcolor=blue,
}

% -- Datos
\institucion{Facultad de Ciencias Exactas, Naturales y Ambientales}
\carrera{Catálogo STEM}
\asignatura{Álgebra lineal}
\tema{Clase 08: Coordenadas y cambio de base}
\autor{Andrés Merino}
\fecha{Periodo 2026-01}

\logouno[0.16\textwidth]{Logos/logoPUCE_04_ac}
\definecolor{colortext}{HTML}{1957d1}
\definecolor{colordef}{HTML}{1957d1}
\fuente{montserrat}

\begin{document}

\encabezado 

%%%%%%%%%%%%%%%%%%%%%%%%%%%%%%%%%%%%%%%%
\section*{Resultado de Aprendizaje}
%%%%%%%%%%%%%%%%%%%%%%%%%%%%%%%%%%%%%%%%

%%%%%%%%%%%%%%%%%%%%%%%%%%%%%%%%%%%%%%%%
\subsection*{RdA de la asignatura:}
\begin{itemize}[leftmargin=*]
    \item \textbf{RdA 1:} Comprender los conceptos fundamentales del Álgebra Lineal, incluyendo el estudio de matrices, determinantes, sistemas de ecuaciones lineales y espacios vectoriales, destacando su importancia en el análisis y resolución de problemas matemáticos.
    \item \textbf{RdA 2:} Resolver problemas prácticos y teóricos relacionados con el Álgebra Lineal, utilizando técnicas de cálculo con matrices, determinantes y sistemas de ecuaciones, así como en el análisis de espacios vectoriales.    
    % \item \textbf{RdA 3:} Aplicar los principios y métodos del Álgebra Lineal en una variedad de contextos demostrando comprensión de su relevancia y utilidad en situaciones prácticas y teóricas.
\end{itemize}


%%%%%%%%%%%%%%%%%%%%%%%%%%%%%%%%%%%%%%%%
\subsection*{RdA de la actividad:}
\begin{itemize}[leftmargin=*]
    \item Calcular el vector de componentes de un vector respecto de una base ordenada.
    \item Construir y utilizar matrices de cambio de base entre dos bases ordenadas.
    \item Relacionar la matriz de cambio de base con su inversa para convertir coordenadas en ambos sentidos.
\end{itemize}

%%%%%%%%%%%%%%%%%%%%%%%%%%%%%%%%%%%%%%%%
\section*{Introducción}
%%%%%%%%%%%%%%%%%%%%%%%%%%%%%%%%%%%%%%%%

%%%%%%%%%%%%%%%%%%%%%%%%%%%%%%%%%%%%%%%%
\paragraph{Control de lectura:}
Antes de iniciar la clase, se aplicará el control de lectura del siguiente enlace:
\href{https://gemini.google.com/share/1a9869cf847a}{Control de lectura - Clase 08}.

\paragraph{Pregunta inicial:}
¿Cómo podemos hacer «traducciones» entre diferentes formas de representar un mismo vector en un espacio vectorial? 

%%%%%%%%%%%%%%%%%%%%%%%%%%%%%%%%%%%%%%%%
\section*{Desarrollo}
%%%%%%%%%%%%%%%%%%%%%%%%%%%%%%%%%%%%%%%%

%%%%%%%%%%%%%%%%%%%%%%%%%%%%%%%%%%%%%%%%
\subsection*{Actividad 1: Coordenadas y cambio de base}

La clase inicia con una recuperación breve de base y dimensión. Luego, mediante clase magistral y práctica guiada, se introduce la noción de base ordenada, vector de componentes y matriz de cambio de base. Finalmente, se trabaja la conversión de coordenadas entre bases y la relación
\[
    (P_{S\gets T})^{-1}=P_{T\gets S}.
\]

\paragraph{¿Cómo lo haremos?}
\begin{itemize}[leftmargin=*]
    \item \textbf{Clase magistral:} se presentan las definiciones de base ordenada, vector de componentes y matriz de cambio de base. Se usan ejemplos en $\mathbb{R}^2$, $\mathbb{R}^3$ y espacios de polinomios. Se usa el resumen \href{https://fcena-puce.github.io/AlgebraLineal-05-N0048/01Resumen/Resumen08.pdf}{Resumen08.pdf}.
    \item \textbf{Resolución de ejercicios:} se guiará a los estudiantes en problemas de cálculo de coordenadas y de construcción de matrices de cambio de base para transformar $[v]_T$ en $[v]_S$ y viceversa. Se puede usar el documento \href{https://fcena-puce.github.io/AlgebraLineal-05-N0048/03Extras/Extra04-CambioBase.pdf}{Extra04-CambioBase.pdf}.
    \item \textbf{Recomendaciones de lectura:} 
    \begin{itemize}
        \item \href{https://blog.nekomath.com/algebra-lineal-i-matrices-de-cambio-de-base/}{Matrices de cambio de base}
    \end{itemize}
\end{itemize}

\paragraph{Verificación de aprendizaje:}
\begin{enumerate}[leftmargin=*]
    \item ¿Cómo se calcula $[u]_B$ cuando se conoce una base ordenada $B$?
    \item ¿Qué información codifica la matriz $P_{S\gets T}$?
    \item ¿Por qué $(P_{S\gets T})^{-1}=P_{T\gets S}$?
\end{enumerate}

%%%%%%%%%%%%%%%%%%%%%%%%%%%%%%%%%%%%%%%%
\section*{Cierre}
%%%%%%%%%%%%%%%%%%%%%%%%%%%%%%%%%%%%%%%%

\paragraph{Tarea:}
Resolver del libro \href{https://catalogobiblioteca.puce.edu.ec/cgi-bin/koha/opac-detail.pl?biblionumber=86081}{Fundamentos de álgebra lineal de Larson}: sección 4.7, ejercicios: 5, 7, 9, 17, 22, 24, 39, 41.

\paragraph{Pregunta de investigación:}
\begin{enumerate}[leftmargin=*]
    \item ¿Qué criterios permiten elegir una base que simplifique el cálculo en un problema aplicado?
    \item ¿Cómo se interpreta geométricamente un cambio de base en $\mathbb{R}^2$ y $\mathbb{R}^3$?
    \item ¿En qué contextos de ingeniería y ciencia de datos se usan cambios de base de manera natural?
\end{enumerate}

\paragraph{Para la próxima clase:}  
Leer la sección 5.2 del libro \href{https://catalogobiblioteca.puce.edu.ec/cgi-bin/koha/opac-detail.pl?biblionumber=86081}{Fundamentos de álgebra lineal de Larson}.

\end{document}
